%CHAPTER 1 Introduction 
\chapter{Introduction}\label{chap:intro}
This chapter delineates the concept of Breach and Attack Simulation (BAS) and establishes its essential role in the testing of security defenses. Examines miter Caldera, Infection Monkey, and Atomic Red Team, specifically focusing on their architecture and capabilities. In conclusion, the chapter reviews earlier studies that have compared these tools to establish the existing state of knowledge.
% 1.1 Motivation
\section{Motivation}\label{sec:Motivation}
Cyber threats are growing rapidly in both volume and sophistication. Attacks now originate from both internal and external sources, traversing multiple channels such as social media, email, SMS, and internal systems\cite{verizon2024}. The widespread availability and ease of use of automated attack tools have lowered the entry barrier for malicious actors, making it easier to breach defenses, steal data, and deploy ransomware\cite{enisa2024}. Given this landscape, organizations must continually improve and validate their security postures.

Organizations often rely on a range of defensive measures to protect against these threats. These include human-centric controls like security awareness training for employees and proactive assessments such as penetration testing and Red Team exercises. They also deploy a vast array of technological defenses, including Intrusion Prevention Systems (IPS), Intrusion Detection Systems (IDS), Antivirus software, Network Firewalls, Web Application Firewalls (WAF), Endpoint Detection and Response (EDR), Network Detection and Response (NDR), and Data Loss Prevention (DLP) solutions\cite{nist80053}. However, even with these extensive tools and methodologies, it remains challenging to definitively assess an organization’s true security posture and its readiness for evolving sophisticated attacks\cite{mandiant2023}.

In response to this challenge, many organizations are adopting Breach and Attack Simulation (BAS), a proactive cybersecurity testing approach that automates the simulation of various cyberattacks to evaluate an organization’s defenses\cite{gartner_bas}. BAS tools simulate real-world attacks—such as malware infiltrations and phishing campaigns—and provide continuous visibility into vulnerabilities across endpoints, networks, and applications. Unlike traditional discrete vulnerability assessments and scans, BAS provides a continuous feedback loop, allowing security teams to quickly address weaknesses as they appear and stay ahead of evolving threats.

While commercial BAS tools are highly effective, their high cost often makes them inaccessible to small or budget-limited organizations. Fortunately, several free and open-source BAS tools have emerged, offering a viable alternative for organizations looking to improve their security posture without significant financial investment.

This research aims to evaluate and compare leading open-source BAS tools to provide a comprehensive, evidence-based guide for organizations. Potential candidates were initially identified by querying five independent AI assistants (Gemini, ChatGPT, Copilot, Claude, and DeepAI) to identify prominent open-source BAS tools for 2025. Based on this consensus, MITRE Caldera\cite{caldera}, Infection Monkey\cite{infectionmonkey}, and Atomic Red Team\cite{atomicredteam} were selected for a comparative analysis. Other suggestions, such as OpenBAS, Metasploit, OWASP ZAP, and Metta, were acknowledged as complementary tools but were excluded due to lower consensus or a primary focus on penetration testing and web application security rather than continuous BAS. The AI recommendations served only as an initial filter; final conclusions are derived from standardized, repeatable lab tests conducted in a controlled environment.

%1.2 Problem Statements
\section{Problem Statements}\label{sec:Hypothesis}
The core problems addressed by this research are defined by the following challenges:
\begin{itemize}
    \item \textbf{Financial Accessibility}: Commercial Breach and Attack Simulation (BAS) platforms are cost-prohibitive for small and medium-sized enterprises (SMEs) and other organizations with limited security budgets.
    \item \textbf{Lack of Practical Usability Data and Implementation Overhead}: Existing information on open-source BAS tools primarily focuses on high-level technical features, leaving a critical knowledge gap concerning their practical usability, measurable implementation complexity, and the interpretability of raw simulation results by non-expert Blue Teams.
    \item \textbf{Absence of Systematic, Evidence-Based Comparison}: There is a scholarly void in systematic, quantitative research that compares the real-world effectiveness and practical operational performance of leading open-source BAS solutions against one another using controlled, empirical metrics.
\end{itemize}

%1.3 Objectives
\section{Objectives}\label{sec:Objectives}
The primary objectives of this research are to:
\begin{itemize}
    \item \textbf{Compare Effectiveness:} Compare how well open-source BAS tools perform in the real world by executing identical, scripted attack scenarios within a controlled lab environment.
    
    \item \textbf{Assess Usability:} Evaluate the operational complexity of each tool by recording the step-by-step process required for installation, configuration, and running simulations.
    
    \item \textbf{Measure \& Rank:} Rank each solution using a Weighted Comparison Matrix, evaluating the tool's holistic effectiveness across six core dimensions: Usability \& Deployment, Threat Intelligence Alignment, Scenario Coverage \& Reliability, Operational Performance, Extensibility, and Community Maturity.
\end{itemize}

%1.4 Scope 
\section{Scope}\label{sec:Scope}
This research is focused on the following key areas:

\begin{itemize}
    \item \textbf{Tool Selection:} The study is limited to three main open-source BAS tools:
    \begin{enumerate}
        \item MITRE Caldera
        \item Infection Monkey
        \item Atomic Red Team.
    \end{enumerate}
    
    These were chosen based on consensus from multiple sources and their relevance to continuous security validation.

    \item \textbf{Operational Focus:} This research focuses strictly on offensive (Red Team) capabilities. Even though some of these tools include defensive (Blue Team) features, this study will only test their ability to simulate attacks. The goal is to evaluate how well they mimic adversary behavior, not how they fix or block issues.

    \item \textbf{Evaluation Environment:} All experiments will be conducted in a controlled laboratory environment using Linux-based virtual machines. Linux was selected because it is open source, requires no subscription licenses, and ensures consistency across all test systems.

    \item \textbf{Scope Limitations:}
    \begin{itemize}
        \item \textbf{Specific Attack Techniques:} We will test a specific set of MITRE ATT\&CK techniques, not the whole framework. (Justification: We selected only the techniques that all three tools can actually perform, ensuring a fair comparison.)

        \item \textbf{Usability over Integration:} The study looks at how well the tools work in a lab, not how they fit into a large corporate network. (Justification: This avoids complications with software licenses and complex setups.)

        \item \textbf{No Commercial Tools:} Paid tools like Cymulate or AttackIQ are not tested, though they may be mentioned for comparison.
    \end{itemize}
\end{itemize}

%1.5 Structure of the Document
\section{Structure of the Document}\label{sec:Structure}
The remainder of this document is organized into the following chapters:

\begin{itemize}
    \item \textbf{Chapter 1: Introduction:} Presents the motivation, problem statements, objectives, scope, and overall structure of the research.
    \item \textbf{Chapter 2: Background and Literature Review:} Provides a comprehensive review of the cybersecurity threat landscape, Breach and Attack Simulation (BAS) concepts, an overview of open-source BAS tools, commercial solutions for context, and related comparative studies.
    \item \textbf{Chapter 3: Proposed Work:} Explains the conceptual framework, selection of BAS tools, experimental design, evaluation criteria, and the weighted scoring matrix used in the study.
    \item \textbf{Chapter 4: Implementation:} Describes the deployment of the selected BAS tools in the laboratory environment, including detailed procedures for installation, configuration, and execution of attack simulations.
    \item \textbf{Chapter 5: Experiment and Results:} Presents the results of the comparative evaluation, including quantitative and qualitative analysis of installation complexity, detection effectiveness, performance, reporting quality, customization capability, and community support.
    \item \textbf{Chapter 6: Conclusion and Future Work:} Summarizes the key findings, highlights the practical implications of the research for budget-conscious organizations, and suggests specific areas for further study and improvement in BAS.
\end{itemize}